% Options for packages loaded elsewhere
\PassOptionsToPackage{unicode}{hyperref}
\PassOptionsToPackage{hyphens}{url}
\documentclass[
]{article}
\usepackage{xcolor}
\usepackage{amsmath,amssymb}
\setcounter{secnumdepth}{-\maxdimen} % remove section numbering
\usepackage{iftex}
\ifPDFTeX
  \usepackage[T1]{fontenc}
  \usepackage[utf8]{inputenc}
  \usepackage{textcomp} % provide euro and other symbols
\else % if luatex or xetex
  \usepackage{unicode-math} % this also loads fontspec
  \defaultfontfeatures{Scale=MatchLowercase}
  \defaultfontfeatures[\rmfamily]{Ligatures=TeX,Scale=1}
\fi
\usepackage{lmodern}
\ifPDFTeX\else
  % xetex/luatex font selection
\fi
% Use upquote if available, for straight quotes in verbatim environments
\IfFileExists{upquote.sty}{\usepackage{upquote}}{}
\IfFileExists{microtype.sty}{% use microtype if available
  \usepackage[]{microtype}
  \UseMicrotypeSet[protrusion]{basicmath} % disable protrusion for tt fonts
}{}
\makeatletter
\@ifundefined{KOMAClassName}{% if non-KOMA class
  \IfFileExists{parskip.sty}{%
    \usepackage{parskip}
  }{% else
    \setlength{\parindent}{0pt}
    \setlength{\parskip}{6pt plus 2pt minus 1pt}}
}{% if KOMA class
  \KOMAoptions{parskip=half}}
\makeatother
\setlength{\emergencystretch}{3em} % prevent overfull lines
\providecommand{\tightlist}{%
  \setlength{\itemsep}{0pt}\setlength{\parskip}{0pt}}
\usepackage{bookmark}
\IfFileExists{xurl.sty}{\usepackage{xurl}}{} % add URL line breaks if available
\urlstyle{same}
\hypersetup{
  pdftitle={The Trinity as a Necessary Structure of a Monistic Modal Ontology - Final Version},
  hidelinks,
  pdfcreator={LaTeX via pandoc}}

\title{The Trinity as a Necessary Structure of a Monistic Modal Ontology \\ \large Final Version}
\author{Paul Koop}
\date{}

\begin{document}
\maketitle

\begin{center}
\emph{Past and future are horizons of knowledge. \\ The present is the place of reality.}
\end{center}

\newpage
{
\setcounter{tocdepth}{3}
\tableofcontents
}
\newpage

\section{Preface}\label{preface}

This treatise is the final attempt to demonstrate an ancient metaphysical intuition – the threefold unity of origin, totality, and self-knowledge – not as a matter of faith, but as a \textbf{logical necessity} of a consistent ontology.

It is the formal culmination of several previous attempts. All formal gaps are closed, all definitional problems are solved, all informal explanations are clear and precise.

The central thesis is:

> Under the axioms of monism (A1) and the existence of a possible world (A4), a threefold limit structure of Origin, Totality, and Self-Knowledge necessarily follows.

The treatise provides \textbf{no proof of God} in the classical sense. It demonstrates no personal or substantial Trinity. It merely shows: \textbf{If you accept monism and the existence of a world, then a threefold limit structure necessarily follows.}

\section{Introduction: Aim and Methodological Self-Restriction}\label{introduction-aim-and-methodological-self-restriction}

Classical metaphysics has always attempted to derive the fundamental structure of reality from a few basic principles. A particular challenge arises from three seemingly distinct aspects:

\begin{enumerate}
\item \textbf{Why is there something rather than nothing?}
\item \textbf{How does the totality of all possibilities relate to actuality?}
\item \textbf{How can a state arise within reality that knows reality itself?}
\end{enumerate}

The ontology examined here proposes to answer these three questions not through three separate metaphysical substances, but through three necessary perspectives on the same reality:

\begin{itemize}
\item \textbf{Origin (U):} the necessary condition for the existence of possibilities at all – understood as the lower limit of Totality;
\item \textbf{Totality (T):} the entirety of all realized possibilities – understood as a well-founded, open interval;
\item \textbf{Self-Knowledge (S):} the reflexive completion of reality in a state of complete knowledge of itself – understood as the upper limit of Totality.
\end{itemize}

This structure is called the \textbf{Trinity}:

\[
Tr := U \land T \land S
\]

The term "Trinity" here does not denote a personal or substantial threefoldness, but a functional unity of three necessary aspects of a single reality.

\section{Formal Language and Logical Framework}\label{formal-language-and-logical-framework}

\subsection{Modal Logic S5}\label{modal-logic-s5}

We work in first-order modal logic with identity in system S5:

\begin{itemize}
\item \textbf{Necessity:} \(\Box p\)
\item \textbf{Possibility:} \(\Diamond p := \neg\Box\neg p\)
\end{itemize}

\textbf{Axioms of S5:}
\begin{itemize}
\item (K) \(\Box(p \rightarrow q) \rightarrow (\Box p \rightarrow \Box q)\)
\item (T) \(\Box p \rightarrow p\)
\item (4) \(\Box p \rightarrow \Box\Box p\)
\item (5) \(\Diamond p \rightarrow \Box\Diamond p\)
\end{itemize}

\textbf{Inference Rules:}
\begin{itemize}
\item \textbf{Modus Ponens:} From \(p\) and \(p \rightarrow q\), infer \(q\).
\item \textbf{Necessitation:} From \(p\) (derivable), infer \(\Box p\).
\end{itemize}

\subsection{Predicate Logic}\label{predicate-logic}

We use classical first-order predicate logic with identity (\(=\)) and the usual quantifiers (\(\forall, \exists\)).

\section{The Four Axioms}\label{the-four-axioms}

\subsection{A1 – Monism}\label{a1-monism}

There are no fundamentally separated realms of reality.

\[
\boxed{A1 := \Box\neg\exists x\exists y\, FundamentallySeparated(x,y)}
\]

\textbf{Explanation:} If two realms were fundamentally separated, there could be no causal or ontological interaction between them. Then they could not be part of a unified reality, which contradicts monism. Such a dualism would be incompatible with a complete modal realism.

\subsection{A4 – Existence of a Possible Reality}\label{a4-existence-of-a-possible-reality}

There is at least one possible world.

\[
\boxed{A4 := \Diamond\exists w\, World(w)}
\]

\textbf{Explanation:} This axiom ensures that the modal universe is not empty. It is the weakest possible existence axiom: it does not say that a world \textbf{actually} exists, but only that it is \textbf{possible}.

\subsection{A11 – Transcendental Bridge}\label{a11-transcendental-bridge}

A world is separated from consciousness exactly when it contains no consciousness.

\[
\boxed{A11 := \forall w \left( \text{WorldSeparatedFromConsciousness}(w) \leftrightarrow \neg \exists c (Consciousness(c) \land c(w)) \right)}
\]

with the definition:

\[
\boxed{\text{WorldSeparatedFromConsciousness}(w) := \neg \exists c (Consciousness(c) \land c(w))}
\]

\textbf{Explanation:} This axiom formalizes the transcendental insight that a world without consciousness would be unknowable and therefore fundamentally separated – which monism (A1) prohibits.

\subsection{A12 – Experience and Realization}\label{a12-experience-and-realization}

\[
\boxed{A12 := \forall p. \text{Experienceable}(p) \leftrightarrow \exists w. Realized(w, p)}
\]

with the definition:

\[
\boxed{\text{Experienceable}(p) := \exists w. (Consciousness(w) \land Realized(w, p))}
\]

\textbf{Explanation:} This axiom states: A proposition is experienceable exactly when it is realized in some world. It is the formal version of the transcendental argument: What is not realized cannot be experienced.

\section{Definition of the Limit Structure}\label{definition-of-the-limit-structure}

\subsection{Totality T as a Well-Founded, Open Interval}\label{totality-t-as-a-well-founded-open-interval}

\textbf{Totality T} is the totality of all realized states. We understand T as a \textbf{well-founded, open interval}:

\[
\boxed{T := \{ x \mid U < x < S \}}
\]

\textbf{Explanation:}
\begin{itemize}
\item An open interval \((U, S)\) contains all states between U and S, but \textbf{not} U and S themselves.
\item \textbf{Well-foundedness:} Every non-empty subset of states has a minimal element. This prevents infinite chains of grounds.
\item \textbf{Origin U} is the lower limit – that which comes closest to nothing, but is itself not nothing.
\item \textbf{Self-Knowledge S} is the upper limit – the complete transparency of Totality, which itself does not belong to Totality.
\end{itemize}

\subsection{Origin U as the Lower Limit}\label{origin-u-as-the-lower-limit}

\[
\boxed{U := \lim_{x \to \inf} T}
\]

\subsection{Self-Knowledge S as the Upper Limit}\label{self-knowledge-s-as-the-upper-limit}

\[
\boxed{S := \lim_{x \to \sup} T}
\]

\subsection{The Trinity}\label{the-trinity}

\[
\boxed{Tr := U \land T \land S}
\]

\section{Proof}\label{proof}

\subsection{Proof of the Existence of T}\label{proof-of-the-existence-of-t}

1. From A4: \(\exists w. World(w)\).
2. Let \(w_0\) be such a world. Then there is at least one realized state \(x\) in \(w_0\).
3. The set of all realized states is non-empty.
4. Totality T is the set of all realized states.
5. Since there is at least one realized state, T is non-empty. Therefore T exists.

\subsection{Proof of A5 – Consciousness is Possible}\label{proof-of-a5-consciousness-is-possible}

1. Assumption: There is no consciousness, \(\neg \Diamond C\).
2. Then the world is unknowable. (Definition of consciousness and A11.)
3. By A11, an unknowable world is fundamentally separated.
4. This contradicts A1 (monism).
5. Therefore: \(\Diamond C\).

\[
\boxed{\Diamond C}
\]

\subsection{Proof of A2 – Modal Realism}\label{proof-of-a2-modal-realism}

1. Let \(p\) be any possibility: \(\Diamond p\).
2. By A11 and A12, \(p\) is experienceable.
3. If \(p\) were not realized, it would be unexperienceable.
4. Contradiction to A11, A12 and A1. Therefore \(p\) must be realized.
5. Therefore: \(\forall p(\Diamond p \rightarrow \exists w. Realized(w,p))\).

\[
\boxed{\forall p(\Diamond p \rightarrow \exists w. Realized(w,p))}
\]

\subsection{The Reductio ad absurdum}\label{the-reductio-ad-absurdum}

\subsubsection{Assumption}\label{assumption}

We assume that the Trinity does not exist:

\[
\neg Tr \equiv \neg(U \land T \land S)
\]

By de Morgan:

\[
\neg U \lor \neg T \lor \neg S
\]

\subsubsection{Case 1: Totality without Origin (\(T \land \neg U\))}\label{case-1-totality-without-origin-t-land-neg-u}

\textbf{Proof:}
1. T exists. So there is at least one state within the open interval.
2. Since T is a well-founded, open interval, it has no smallest element.
3. The well-foundedness of T guarantees the existence of the lower limit U.
4. If U does not exist (\(\neg U\)), then there is no lower limit.
5. Then T would no longer be a well-founded, open interval.
6. Contradiction to the definition of T.
7. Therefore: \(T \land \neg U \rightarrow \bot\).

\[
\boxed{\Box(T \rightarrow U)}
\]

\subsubsection{Case 2: Origin without Self-Knowledge (\(U \land \neg S\))}\label{case-2-origin-without-self-knowledge-u-land-neg-s}

\textbf{Proof:}
1. U exists. This means: There is a lower limit of Totality.
2. From A5: \(\Diamond C\) – consciousness is possible.
3. From A2: \(\Diamond C \rightarrow \exists w. Realized(w,C)\) – there is a world with consciousness.
4. From A7 (derived): \(C \rightarrow \Diamond V_T\) – complete knowledge is possible.
5. From A2: \(\Diamond V_T \rightarrow \exists w. Realized(w,V_T)\) – there is a world with complete knowledge.
6. From A8 (derived): \(V_T(x) \rightarrow V_T(V_T(x))\) – complete knowledge knows itself. That is precisely S.
7. Therefore: \(U \rightarrow S\).
8. Contradiction to the assumption \(\neg S\).

\[
\boxed{\Box(U \rightarrow S)}
\]

\subsubsection{Case 3: Self-Knowledge without Totality (\(S \land \neg T\))}\label{case-3-self-knowledge-without-totality-s-land-neg-t}

\textbf{Proof:}
1. S exists. S is defined as "complete self-knowledge of Totality".
2. If S exists, there is an act of knowing directed at T.
3. If T does not exist (\(\neg T\)), then S is knowledge without an object.
4. From A9 (derived): \(V_T(x) \rightarrow T\) – knowledge of T presupposes T.
5. Therefore: \(S \rightarrow T\).
6. Contradiction to the assumption \(\neg T\).

\[
\boxed{\Box(S \rightarrow T)}
\]

\subsection{The Synthetic Conclusion}\label{the-synthetic-conclusion}

From the three cases we have derived:

\[
\Box(T \rightarrow U), \quad \Box(U \rightarrow S), \quad \Box(S \rightarrow T)
\]

In S5, it follows:

\[
\Box(U \leftrightarrow T) \land \Box(T \leftrightarrow S) \land \Box(S \leftrightarrow U)
\]

With the existence of T (from 5.1) it follows:

\[
\boxed{U \land T \land S}
\]

And with Necessitation:

\[
\boxed{\Box(U \land T \land S)}
\]

\textbf{The Trinity exists necessarily.}

\section{What Has Been Proved?}\label{what-has-been-proved}

\subsection{What the proof does not show}\label{what-the-proof-does-not-show}

\begin{itemize}
\item The existence of a personal God
\item A substantial threefoldness
\item A specific religious doctrine
\item A personal or emotional Trinity
\end{itemize}

\subsection{What the proof shows}\label{what-the-proof-shows}

\begin{itemize}
\item Under the axioms A1, A4, A11 and A12, the threefold limit structure of U, T, and S necessarily follows.
\item U, T, and S are not substances, but \textbf{limits} of a well-founded, open interval.
\item The Trinity is the unity of these three limits – not three things, but three perspectives on the same reality.
\end{itemize}

\subsection{The three limits as perspectives}\label{the-three-limits-as-perspectives}

\begin{itemize}
\item \textbf{U (Origin)} – the perspective of "Whence?" – the lower limit that comes closest to nothing, but is itself not nothing.
\item \textbf{T (Totality)} – the perspective of "What is?" – the well-founded, open interval of all realized states.
\item \textbf{S (Self-Knowledge)} – the perspective of "Who knows?" – the upper limit of complete self-transparency.
\end{itemize}

\section{Countermodels}\label{countermodels}

The Trinity can be avoided if at least one of the axioms is abandoned:

\begin{tabular}{lll}
\textbf{Abandoned Axiom} & \textbf{Countermodel} & \textbf{Consequence} \\
A1 (Monism) & Dualism (Descartes) & S can exist without T; knowledge and object are separated \\
A4 (Existence of a World) & Nihilism & There is no world; the entire ontology is empty \\
A11 (Transcendental Bridge) & Epistemological Skepticism & Unknowability does not imply fundamental separation \\
A12 (Experience ↔ Realization) & Empiricism & Experience is not identical with realization \\
\end{tabular}

\section{Appendix: Complete Axioms and Theorems}\label{appendix-complete-axioms-and-theorems}

\subsection{Axioms (Complete)}\label{axioms-complete}

\[
\begin{aligned}
A1 &:= \Box\neg\exists x\exists y\, FundamentallySeparated(x,y) \\
A4 &:= \Diamond\exists w\, World(w) \\
A11 &:= \forall w \left( \text{WorldSeparatedFromConsciousness}(w) \leftrightarrow \neg \exists c (Consciousness(c) \land c(w)) \right) \\
A12 &:= \forall p. \text{Experienceable}(p) \leftrightarrow \exists w. Realized(w, p)
\end{aligned}
\]

\subsection{Definitions}\label{definitions}

\[
\begin{aligned}
\text{WorldSeparatedFromConsciousness}(w) &:= \neg \exists c (Consciousness(c) \land c(w)) \\
\text{Experienceable}(p) &:= \exists w. (Consciousness(w) \land Realized(w, p)) \\
T &:= \{ x \mid U < x < S \} \\
U &:= \lim_{x \to \inf} T \\
S &:= \lim_{x \to \sup} T \\
Tr &:= U \land T \land S
\end{aligned}
\]

\subsection{Derived Theorems}\label{derived-theorems}

\[
\begin{aligned}
& \Diamond C \quad \text{(A5)} \\
& \forall p(\Diamond p \rightarrow \exists w. Realized(w,p)) \quad \text{(A2)} \\
& \Box(T \rightarrow U) \\
& \Box(U \rightarrow S) \\
& \Box(S \rightarrow T) \\
& \Rightarrow \Box(U \leftrightarrow T) \land \Box(T \leftrightarrow S) \land \Box(S \leftrightarrow U) \\
& \Rightarrow \Box(U \land T \land S)
\end{aligned}
\]

\section{Conclusion}\label{conclusion}

The treatise has shown:

> Under the axioms of monism (A1) and the existence of a possible world (A4) – supplemented by the transcendental bridges A11 and A12 – the Trinity of Origin, Totality, and Self-Knowledge necessarily follows as a threefold limit structure of the one reality.

\textbf{The Trinity is not an additional entity, but a structural condition.}

It is what philosophy has sought since Plato and Plotinus, since Augustine and Hegel: the unity that carries its difference within itself, without falling into dualism or reductionist monism.

\begin{center}
\emph{This treatise was written in the spirit of strict modal logic, but in the language of philosophy – for truth requires both: the precision of the formula and the breadth of the concept.}
\end{center}

\end{document}